\chapter{Conclusion and Future Work}


\section{Overall Summary}

This thesis addresses the core issue of improving grading efficiency and feedback quality in programming courses by designing and implementing an AI-assisted teaching system. Unlike approaches focused on isolated algorithms, this project emphasizes an end-to-end engineering workflow in teaching contexts, covering submission, controlled execution, static analysis, scoring feedback, instructor management, and deployment operation.

At implementation level, this work completes backend services, database modeling, API partitioning, and frontend-backend integration, and establishes a reproducible containerized deployment scheme with off-campus access support. Through instructor-side class and student detail functions, the system supports not only scoring but also traceable management views. Test results indicate good usability and stability under course-scale scenarios.

At method level, final scores are separated from static-analysis scores, and template/blank-submission detection is integrated to balance correctness and process evaluation. Meanwhile, RAG is introduced to align feedback with course context and reduce generic or out-of-scope suggestions. Overall, this work shows that integrated deployment of "automatic grading + rule constraints + retrieval-enhanced feedback + teaching management" is feasible in undergraduate teaching settings.

\section{Limitations}

Although the system already supports daily teaching demos and small-scale use, several limitations remain. First, feedback quality is affected by model service status and retrieval hit quality, and may still be insufficiently specific in edge cases. Second, current performance evaluation focuses mainly on course-scale integration and does not yet cover higher concurrency and longer continuous operation. Third, instructor-side analytics remain relatively basic, with room to improve metric granularity and visualization depth.

In addition, this thesis focuses mainly on backend and database implementation. Experiments on frontend interaction quality and student learning behavior are not yet sufficiently deep. For example, there is still limited systematic comparative evidence on how different feedback presentation styles affect learning outcomes. These issues provide clear directions for future work.

\section{Future Work}

Future development can proceed in three directions. First, continue optimizing the feedback pipeline by expanding course corpora, improving retrieval ranking strategies, and strengthening feedback template constraints to improve output stability and teaching alignment. Second, strengthen engineering robustness by introducing automated regression and stress tests, optimizing long-latency task handling, and improving failure-recovery mechanisms. Third, enrich teaching analytics by adding error-trend statistics, class-level learning trajectory observation, and finer-grained instructor decision-support views.

From an application perspective, the system can be extended to more assignment types and programming languages, and further integrated with course platforms. With iterative updates based on real classroom data, the system has the potential to evolve from an "auxiliary grading tool" into a "teaching decision-support platform component," improving both teaching efficiency and student feedback experience.

